\subsection{模的加性函数}

设 A 是一个环，C 是 A-模的类所成的一个集（VIII, 第 51 页）；若一个 A-模的类属于 C，则称它是 C 型的。

定义 1. —— 若每个零模都是 C 型的，且两个 C 型模的直和是 C 型的，则称 A-模的类所成之集 C 是加性的。若 C 是加性的，且一个 C 型模的子模与商模都是 C 型的，则称 C 是遗传的。

例. —— 1) 有限长度的 A-模的类所成之集是遗传的（II, §1, No.~10, 第 212 页，命题 16）。

\begin{enumerate}
\def\labelenumi{\arabic{enumi})}
\setcounter{enumi}{1}
\item
  有限生成的 A-模的类所成之集是加性的。若环 A 是诺特的，则这个集是遗传的（VIII, 第 3 页，命题 3 与 VIII, 第 7 页，命题 4）。
\item
  有限生成的投射 A-模的类所成之集是加性的，但一般不是遗传的。
\end{enumerate}

定义 2. —— 设 \(\varphi\) 是从 \(\mathcal{C}\) 到一个交换群 G（写作加法）中的一个映射；令 \(\varphi(\mathrm{E}) = \varphi(\mathrm{cl}(\mathrm{E}))\)（对每个 A-模 E，\(\mathcal{C}\) 型的）。我们称 \(\varphi\) 是模的一个加性函数（相应地，弱加性函数），如果对每个正合列（相应地，对每个分裂正合列）都有 \(\varphi(\mathrm{E}) = \varphi(\mathrm{E}') + \varphi(\mathrm{E}'')\)，该列由

\[
0 \longrightarrow \mathrm{E} ^ {\prime} \longrightarrow \mathrm{E} \longrightarrow \mathrm{E} ^ {\prime \prime} \longrightarrow 0
\]

C 型模组成。

例 4. —— 设 C 是有限长度的 A-模的类所成之集。将有限长度的 A-模的类映到其长度的映射 \(\operatorname{long}_{\mathbf{A}} : \mathcal{C} \to \mathbf{Z}\) 是模的一个加性函数（II, §1, No.~10, 第 213 页，推论 3）。本小节的结果是 II, §1, No.~10, 第 212--214 页中关于有限长度模所建立结果的一个推广。

在本小节的余下部分，我们考虑 A-模的一个加性集 C 与一个从 C 到交换群 G 的加性映射 \(\varphi\)。

设 \(E\) 与 \(E'\) 是 \(\mathcal{C}\) 型模；则 \(E \oplus E'\) 是 \(\mathcal{C}\) 型的，且存在一个分裂正合列（II, §1, No.~9, 第 210 页）

\[
0 \longrightarrow \mathrm{E} \longrightarrow \mathrm{E} \oplus \mathrm{E} ^ {\prime} \longrightarrow \mathrm{E} ^ {\prime} \longrightarrow 0;
\]

由此我们得出

\[
\varphi (\mathrm{E} \oplus \mathrm{E} ^ {\prime}) = \varphi (\mathrm{E}) + \varphi (\mathrm{E} ^ {\prime}).\tag{1}
\]

特别地，我们有 \(\varphi(0) = 0\)。

命题 1. —— 假设 \(\mathcal{C}\) 是遗传的。设 E 与 F 是 A-模，\(u: \mathrm{E} \to \mathrm{F}\) 是一个线性映射。

\begin{enumerate}
\def\labelenumi{\alph{enumi})}
\item
  若 E 或 F 是 C 型的，则 u 的象也是 C 型的。
\item
  若 E 是 C 型的，则 u 的核也是 C 型的，且我们有
\end{enumerate}

\[
\varphi (\mathrm{E}) = \varphi (\mathrm{Ker} u) + \varphi (\mathrm{Im} u).\tag{2}
\]

\begin{enumerate}
\def\labelenumi{\alph{enumi})}
\setcounter{enumi}{2}
\tightlist
\item
  若 \(\mathrm{F}\) 是 \(\mathcal{C}\) 型的，则 \(u\) 的余核亦然，且我们有
\end{enumerate}

\[
\varphi (\mathrm{F}) = \varphi (\operatorname{Im} u) + \varphi (\operatorname{Coker} u).\tag{3}
\]

该命题由如下正合列的存在性得出：

\[
0 \longrightarrow \operatorname{Ker} u \longrightarrow \mathrm{E} \longrightarrow \operatorname{Im} u \longrightarrow 0,
\]

\[
0 \longrightarrow \operatorname{Im} u \longrightarrow \mathrm{F} \longrightarrow \operatorname{Coker} u \longrightarrow 0.
\]

推论. —— 设 \((\mathrm{E}_{i})_{0\leqslant i\leqslant n}\) 是 C 型模的一个有限序列。若存在正合列

\[
0 \longrightarrow \mathrm{E} _ {0} \stackrel {u _ {0}} {\longrightarrow} \mathrm{E} _ {1} \stackrel {u _ {1}} {\longrightarrow} \dots \longrightarrow \mathrm{E} _ {n - 1} \stackrel {u _ {n - 1}} {\longrightarrow} \mathrm{E} _ {n} \longrightarrow 0,
\]

则我们有

\[
\sum_ {i = 0} ^ {n} (- 1) ^ {i} \varphi (\mathrm{E} _ {i}) = 0.\tag{4}
\]

我们对 \(n\) 用归纳法证明该推论；情形 \(n = 0\) 与 \(n = 1\) 是平凡的。设 \(n \geqslant 1\)，且设

\[
0 \longrightarrow \mathrm{E} _ {0} \stackrel {u _ {0}} {\longrightarrow} \mathrm{E} _ {1} \stackrel {u _ {1}} {\longrightarrow} \dots \longrightarrow \mathrm{E} _ {n - 1} \stackrel {u _ {n - 1}} {\longrightarrow} \mathrm{E} _ {n} \stackrel {u _ {n}} {\longrightarrow} \mathrm{E} _ {n + 1} \longrightarrow 0
\]

是一个 \(\mathcal{C}\) 型模的正合列。由命题 1，\(F\)（即 \(u_{n}\) 的核）是一个 \(\mathcal{C}\) 型模，且我们有

\[
\varphi (\mathrm{F}) = \varphi (\mathrm{E} _ {n}) - \varphi (\mathrm{E} _ {n + 1}).\tag{5}
\]

此外，我们有正合列

\[
0 \longrightarrow \mathrm{E} _ {0} \longrightarrow \mathrm{E} _ {1} \longrightarrow \mathrm{E} _ {2} \longrightarrow \dots \longrightarrow \mathrm{E} _ {n - 1} \longrightarrow \mathrm{F} \longrightarrow 0,
\]

且归纳假设给出关系式

\[
\sum_ {i = 0} ^ {n - 1} (- 1) ^ {i} \varphi (\mathrm{E} _ {i}) + (- 1) ^ {n} \varphi (\mathrm{F}) = 0.\tag{6}
\]

于是由 (5) 与 (6) 立即得出 \(\sum_{i=0}^{n+1}(-1)^{i}\varphi(\mathrm{E}_{i}) = 0\)，推论得证。

命题 2. —— 假设集 C 是遗传的。设 E 是一个 A-模，M 与 N 是 E 的子模。

\begin{enumerate}
\def\labelenumi{\alph{enumi})}
\tightlist
\item
  若模 M 与 N 是 C 型的，则模 \(M \cap N\) 与 \(M + N\) 也是 C 型的，且我们有
\end{enumerate}

\[
\varphi (\mathrm{M} + \mathrm{N}) + \varphi (\mathrm{M} \cap \mathrm{N}) = \varphi (\mathrm{M}) + \varphi (\mathrm{N}).
\]

\begin{enumerate}
\def\labelenumi{\alph{enumi})}
\setcounter{enumi}{1}
\tightlist
\item
  若模 E/M 与 E/N 是 C 型的，则模 \(\mathrm{E}/(\mathrm{M}\cap\mathrm{N})\) 与 \(\mathrm{E}/(\mathrm{M}+\mathrm{N})\) 也是 C 型的，且我们有
\end{enumerate}

\[
\varphi (\mathrm{E} / (\mathrm{M} + \mathrm{N})) + \varphi (\mathrm{E} / (\mathrm{M} \cap \mathrm{N})) = \varphi (\mathrm{E} / \mathrm{M}) + \varphi (\mathrm{E} / \mathrm{N}).
\]

断言 a) 与 b) 由如下正合列的存在性得出：

\[
0 \to \mathrm{M} \cap \mathrm{N} \to \mathrm{M} \oplus \mathrm{N} \to \mathrm{M} + \mathrm{N} \to 0
\]

与

\[
0 \to \mathrm{E} / (\mathrm{M} \cap \mathrm{N}) \to (\mathrm{E} / \mathrm{M}) \oplus (\mathrm{E} / \mathrm{N}) \to \mathrm{E} / (\mathrm{M} + \mathrm{N}) \to 0
\]

（II, §1, No.~7, 第 207 页，命题 10）。

命题 3. —— 假设 \(\mathcal{C}\) 是遗传的。设 E 是一个 \(\mathcal{C}\) 型模，\((\mathrm{E}_i)_{0\leqslant i\leqslant n}\) 是 E 的一个合成列（I, §4, No.~7, 第 41 页）。对 \(1\leqslant i\leqslant n\)，模 \(\mathrm{E}_{i - 1} / \mathrm{E}_i\) 是 \(\mathcal{C}\) 型的，且我们有

\[
\varphi (\mathrm{E}) = \sum_ {i = 1} ^ {n} \varphi (\mathrm{E} _ {i - 1} / \mathrm{E} _ {i}).
\]

由于 C 是遗传的，模 \(E_{0}/E_{1}\) 与 \(E_{1}\) 是 C 型的，且我们有 \(\varphi(\mathrm{E}) = \varphi(\mathrm{E}_{0}/\mathrm{E}_{1}) + \varphi(\mathrm{E}_{1})\)。由于序列 \((\mathrm{E}_{i+1})_{0 \leqslant i \leqslant n-1}\) 是 \(E_{1}\) 的一个合成列，命题 3 由对 n 的归纳法得出。

